Introduïm una espira metàl·lica rectangular de $5\ \Omega$ de resistència elèctrica en una regió de l'espai delimitada per un camp magnètic uniforme de 0,2~T perpendicular a la superfície de l'espira. Les dimensions de l'espira són $a = 3\ \mathrm{cm}$ i $b = 6\ \mathrm{cm}$, i es mou a una velocitat de $2\ \mathrm{m\,s^{-1}}$.

\figura[0.8]{fig1.png}

\begin{apartat}{1}
Digueu si circula corrent elèctric per l'espira en les tres situacions següents: en entrar al camp, quan hi està totalment immersa i en sortir-ne, i determineu en cada cas el sentit de circulació de la intensitat corresponent. Justifiqueu les respostes.
\end{apartat}

\begin{apartat}{1}
Calculeu la força electromotriu i la intensitat del corrent elèctric que es genera en cada cas.
\end{apartat}
